\subsubsection{Klasifikasi Ancaman}

Klasifikasi ancaman dilakukan menggunakan pendekatan \textit{STRIDE-per-Interaction}. Setiap interaksi antar elemen pada Data Flow Diagram (Gambar~\ref{fig:decmed-dfd}) diperiksa berdasarkan arah dan tipe interaksinya (\textit{process} ke \textit{data store}, \textit{process} ke \textit{process}, \textit{data flow} yang melintasi \textit{machine boundary}, dsb.), kemudian kategori STRIDE yang dianalisis dibatasi hanya pada kategori yang relevan bagi tipe interaksi tersebut, mengacu pada Tabel~\ref{tab:stride-per-interaction}. Pendekatan ini memberikan justifikasi eksplisit atas kategori yang dipilih maupun yang tidak dianalisis untuk setiap ancaman.

Elemen \texttt{E6 IOTA Network} dan \texttt{E7 IOTA Gas Station} pada \texttt{TB5 External Provider} diperlakukan sebagai infrastruktur eksternal tepercaya yang berada di luar kendali sistem DecMed; ancaman terhadap operasional internal keduanya berada di luar cakupan \textit{threat model} ini, dan hanya interaksinya dengan sistem (melalui P8, sebagai \textit{process} yang mengirim output ke \textit{external interactor}) yang dianalisis.

Hasil klasifikasi ancaman disajikan pada Tabel~\ref{tab:klasifikasi-ancaman}, dengan kolom \textbf{Tipe Interaksi} merujuk pada nomor baris interaksi di Tabel~\ref{tab:stride-per-interaction} yang menjadi dasar penentuan kategori STRIDE bagi ancaman terkait.

\begin{xltabular}{\textwidth}{|>{\centering\arraybackslash}p{0.8cm}|X|>{\centering\arraybackslash}p{1.3cm}|>{\centering\arraybackslash}p{1.5cm}|>{\centering\arraybackslash}p{1.5cm}|>{\centering\arraybackslash}p{1.5cm}|}
    \caption{Klasifikasi Ancaman Sistem DecMed} 
    \label{tab:klasifikasi-ancaman} \\
    
    \hline
    \textbf{ID} & \textbf{Deskripsi Ancaman} & \textbf{Elemen DFD} & \textbf{Level Kepercayaan} & \textbf{Tipe Interaksi} & \textbf{STRIDE} \\
    \hline
    \endfirsthead

    \hline
    \textbf{ID} & \textbf{Deskripsi Ancaman} & \textbf{Elemen DFD} & \textbf{Level Kepercayaan} & \textbf{Tipe Interaksi} & \textbf{STRIDE} \\
    \hline
    \endhead

    \hline
    \endfoot

    \hline
    \endlastfoot

    T1 & Penyerang mendapatkan kredensial atau membajak sesi lokal milik pengguna untuk berpura-pura menjadi pemilik identitas sah saat \textit{sign-in}. & E1--E5, P1, P2, P3 & LK1--LK6 & 7, 11 & S \\
    \hline

    T2 & Kredensial \textit{sign-in} (PIN) tersadap atau terekam melalui \textit{shoulder surfing}, \textit{keylogger}, atau malware pada perangkat saat dimasukkan ke Client. & E1--E5, P1, P2, P3 & LK1--LK6 & 7, 11 & I \\
    \hline

    T3 & Percobaan \textit{sign-in} berulang secara masif menghabiskan sumber daya Client (\textit{local resource exhaustion}). & E1--E5, P1, P2, P3 & LK1 & 7 & D \\
    \hline

    T4 & Payload QR Access Delegation dipalsukan sebelum dipindai oleh Patient Client. & P2, F2 & LK2, LK4 & 7 & S \\
    \hline

    T5 & Payload QR Access Delegation disadap sehingga informasi delegasi akses bocor kepada pihak tidak berwenang. & P2, F2 & LK2, LK4 & 7 & I \\
    \hline

    T6 & Pasien menyangkal pernah memberikan akses (\texttt{create\_access}) kepada tenaga medis tertentu. & E1, P2, P8 & LK2 & 11 & R \\
    \hline

    T7 & Tenaga medis atau tenaga administratif menyangkal pernah mengakses maupun mengubah data pasien tertentu. & E2, E3, P1, P6 & LK3, LK4 & 11 & R \\
    \hline

    T8 & Admin fasyankes menyangkal pernah menerbitkan \textit{activation key} kepada personel tertentu. & E4, P1 & LK5 & 11 & R \\
    \hline

    T9 & Ministry Administrator menyangkal pernah meregistrasi fasyankes tertentu beserta admin-nya. & E5, P3 & LK6 & 11 & R \\
    \hline

    T10 & Penyerang menyamar sebagai Proxy Re-encryption API (P4) yang sah (\textit{man-in-the-middle}) saat Client mengirim permintaan, karena flow melintasi batas perangkat (\textit{client}--\textit{backend}). & P1, P2, P4, F7, F8 & LK2--LK4 & 2, 8 & S \\
    \hline

    T11 & Penyerang menyadap komunikasi Client--PRE Server yang melintasi batas perangkat, mengungkap token JWT maupun data administratif. & P1, P2, P4, F7, F8, F15 & LK2--LK4 & 8 & I \\
    \hline

    T12 & Pihak eksternal membanjiri \textit{endpoint} \texttt{POST /nonce} atau \texttt{POST /keys} dengan permintaan palsu untuk menghabiskan sumber daya PRE Server. & P4, P5, F7, F8 & LK1 & 7 & D \\
    \hline

    T13 & Token JWT dipalsukan atau diputar ulang (\textit{replay}) pada permintaan dari Proxy Re-encryption API ke Medical Record Management. & P4, P6, F12 & LK3, LK4 & 2 & S \\
    \hline

    T14 & Penyadapan jalur \textit{dispatch} internal dari Proxy Re-encryption API ke Authentication/Medical Record Management mengungkap data yang sedang diproses. & P4, P5, P6, F11, F12 & LK7 & 2 & I \\
    \hline

    T15 & Penyerang memodifikasi \textit{Move transaction} pada jalur dari IOTA Transaction Client menuju IOTA Network, karena flow melintasi batas sistem DecMed dengan infrastruktur eksternal. & P8, F23 & LK7, LK8 & 3, 8 & T \\
    \hline

    T16 & Penyerang menyamar sebagai IOTA Transaction Client (P8) yang sah saat berinteraksi dengan Gas Station untuk mengajukan sponsorship transaksi palsu. & P8, F24 & LK7, LK8 & 3 & S \\
    \hline

    T17 & PRE Server (P8) menyangkal telah mengajukan atau menerima keputusan sponsorship transaksi tertentu dari Gas Station. & P8, F24, F25 & LK7, LK8 & 3 & R \\
    \hline

    T18 & Penyerang mengajukan permintaan sponsorship \textit{gas} secara berulang sehingga \textit{gas object} habis dan transaksi sah pengguna lain gagal tersponsori. & P8, F24, F25 & LK1--LK4 & 3 & D \\
    \hline

    T19 & Penyerang dengan akses tidak sah memodifikasi kunci privat atau material PRE yang tersimpan pada Client Key Store. & D3, P1, P2, F5, F6 & LK1 & 5, 10 & T \\
    \hline

    T20 & Penyerang dengan akses fisik ke perangkat yang tidak terkunci membaca kunci privat pada Client Key Store. & D3, F5, F6 & LK1 & 10 & I \\
    \hline

    T21 & Pihak dengan akses tidak sah ke Redis memodifikasi nilai \textit{nonce} autentikasi atau \texttt{kfrag} sebelum digunakan. & D1, P4, P5, P7, F13, F14 & LK7 & 5, 10 & T \\
    \hline

    T22 & Pihak dengan akses tidak sah ke Redis membaca \texttt{kfrag} atau \textit{nonce} sebelum masa berlakunya (TTL) berakhir. & D1 & LK7 & 10 & I \\
    \hline

    T23 & Redis dipenuhi (\textit{memory exhaustion}) melalui permintaan \textit{nonce} baru secara berulang sehingga proses autentikasi pengguna sah terganggu. & D1, F13 & LK1 & 5 & D \\
    \hline

    T24 & Data klinis terenkripsi diunduh dari IPFS Cluster Storage oleh pihak tidak berwenang; risiko pengungkapan tetap ada apabila kunci enkripsi terkait bocor atau lemah. & D2, F18, F19 & LK1, publik & 10 & I \\
    \hline

    T25 & IPFS Cluster dibanjiri permintaan \textit{fetch}/\textit{pin} berukuran besar sehingga mengganggu ketersediaan data klinis bagi pengguna sah. & D2 & publik & 5 & D \\
    \hline

    T26 & Metadata \textit{on-chain} yang tersimpan pada IOTA On-Chain State bersifat publik sehingga pola transaksi berpotensi dianalisis untuk mengungkap informasi tidak langsung mengenai pasien. & D4, F21, F22 & LK9, publik & 10 & I \\
    \hline

    T27 & Kelemahan validasi \textit{capability on-chain} memungkinkan pembacaan atau pembuatan \texttt{ProxyCapEnum}/\texttt{GlobalAdminCap} tanpa otorisasi sah, setara privilese tertinggi PRE Server/Kementerian Kesehatan. & P7, D4 & LK2/LK4 $\rightarrow$ LK6/LK7 & 5 & E \\
    \hline

    T28 & Kegagalan validasi otorisasi pada Authentication \& Access Control memungkinkan Medical Record Management memproses permintaan dengan \textit{scope} lebih tinggi dari yang seharusnya (administratif $\rightarrow$ klinis). & P5, P6, F26 & LK3 $\rightarrow$ LK4 & 2 & E \\
    \hline

\end{xltabular}